\section{Related work}
\subsection{Learning visual speech representations}
Continuous lip reading combines spatial analysis of the mouth with temporal and
linguistic modeling. Sentence-level systems established the importance of learning
over variable-length visual sequences \cite{chung2017lipreading,afouras2022deep}.
Large-scale supervised recognition subsequently expanded the available training
resources \cite{shillingford2019large}. Auto-AVSR uses automatically generated
labels to exploit additional audiovisual collections \cite{ma2023autoavsr}, while
AV-HuBERT learns audiovisual representations through masked cluster prediction
\cite{shi2022avhubert}. These approaches address representation quality and scale.
Our experiments instead start from existing language-specific recognizers and
study how the presentation of supervision affects transfer under a fixed update
budget, using the same visual backbone and source-video collection.

The recognizers considered here use a hybrid CTC/attention formulation
\cite{graves2006ctc,watanabe2017hybrid}. The two output components provide useful
controls because changing their relative decoding weight can change substitutions,
deletions and insertions without changing the trained model. Accordingly, we report
both a decoder selected for each checkpoint on development data and a common
decoder. Together, these comparisons distinguish the performance of the tuned
system from the effect observed under a shared decoding rule.

\subsection{Temporal crops and short-sequence supervision}
Training on shorter sequences is not new to lip reading. Early sentence-level work
used a progression in sequence difficulty \cite{chung2017lipreading}.
Martinez et al. vary temporal context around an intact target word, improving
robustness to frame removal in isolated-word recognition
\cite{martinez2020lipreading}. Their BGRU comparison also illustrates that
robustness and unperturbed accuracy can change in different directions.
Prajwal et al. train with consecutive-word subsequences and compare increasing
subsequence length with using different lengths together
\cite{prajwal2022subword}. Their study is particularly relevant because it shows
that mixed temporal scales provide an established augmentation baseline. We build
on this setting to examine transfer under matched continuation budgets and
whole-utterance evaluation on a different corpus.

Recent word-level work also examines temporal information and augmentation.
InfoSyncNet combines temporal modeling with time masking, mixup and word-boundary
indicators, reporting different time-masking effects on LRW and LRW1000
\cite{xue2025infosyncnet}. These experiments retain an isolated-word recognition
task. FM instead changes both the visual interval and its transcript during
continuation, then evaluates complete utterances on a destination corpus.

The specific setting examined here differs in two respects. The fragment is defined
by a fixed time interval rather than a fixed number of complete words. Its label
contains any transcript unit that overlaps the interval, so boundary words can
remain in the target despite incomplete visual coverage. Moreover, evaluation
uses complete utterances from a different corpus, rather than the same fragment
task used during continuation. These choices create a testable question about the
interaction between context availability and transfer. We examine this fixed-time
formulation through paired context-restoration controls and a direct comparison
with conventional word-aligned fragments. The latter tests the crop choice within
our continuation protocol; it is not a reimplementation of Prajwal et al.'s
architecture or complete training pipeline. The distinction from prior work lies
in the joint assessment of whole-utterance transfer, source performance and
decoder dependence after continuation, rather than in temporal cropping alone.

\subsection{Contextual targets and consistency learning}
Self-distillation provides a way to use information that is available during
training but absent from a particular input view. For example, data2vec predicts
contextual representations of the full input from a masked view
\cite{baevski2022data2vec}. Closer to the recognition objective, CR-CTC encourages
agreement between framewise CTC distributions from two augmented speech views
\cite{yao2025crctc}. Its branches share transcript supervision, with bidirectional
consistency and stop-gradient targets at corresponding time positions.

Our full-view control examines broader temporal context during fragment learning.
The current model supplies a stop-gradient CTC target from the complete video,
restricted to the fragment's time positions, while the student retains its
fragment transcript. It therefore tests contextual guidance alongside the
alternative of exposing additional context directly to the encoder. Both controls
retain ordinary whole-utterance inference. This framing connects the experiments
to established consistency learning while focusing the comparison on the context
available when learning from partial speech.

\subsection{Cross-corpus evaluation and multilingual evidence}
Speaker adaptation addresses distribution shift using examples from the intended
users. Kim et al. learn prompts from target-speaker adaptation data while
preserving the pretrained recognizer parameters \cite{kim2025prompt}. Here,
recognizer gradient updates use the continuation source; destination development
data are used only for decoder calibration where the protocol permits it.
This distinction separates source continuation from target-speaker fitting.

WildVSR was introduced to examine whether strong VSR models generalize beyond
LRS3 \cite{djilali2024wildvsr}. Its results motivate retaining an external-corpus
assessment even when source-domain validation appears satisfactory. We use WildVSR
as the English destination corpus and ViCocktail as the Vietnamese destination
\cite{nguyen2025vicocktail}. The Vietnamese source is the VLR dataset
\cite{vlr2026dataset}; English continuation uses prepared LRS2 training clips.
The two comparisons share the intervention and whole-utterance objective, but
they are distinct transfer settings with different initializations, tokenizers,
training coverage and development-data availability.

We analyze each corpus pair separately, examining training coverage, edit types
and sequence length as possible explanations for their contrasting outcomes.
